\documentclass[11pt]{article}
% ===========================================================================
%  Shared preamble for the AP Calculus AB exam-preparation worksheets.
%  Input right after \documentclass. Files beginning with "_" are not compiled
%  on their own. Navy and gold, to match the course's lesson pages.
% ===========================================================================
\usepackage[letterpaper,top=2.3cm,bottom=2.2cm,left=1.7cm,right=1.7cm,headheight=28pt,footskip=30pt]{geometry}
\usepackage{amsmath,amssymb}
\usepackage[T1]{fontenc}
\usepackage[utf8]{inputenc}
\usepackage{lmodern}
\usepackage{xcolor}
\usepackage[most]{tcolorbox}
\usepackage{enumitem}
\usepackage{fancyhdr}
\usepackage{array}
\usepackage{tikz}
\usepackage{pgfplots}
\pgfplotsset{compat=1.18}

\definecolor{navy}{HTML}{161E54}
\definecolor{ink}{HTML}{27357E}
\definecolor{gold}{HTML}{B98A00}
\definecolor{goldbg}{HTML}{FFFDF3}
\definecolor{inkbg}{HTML}{F6F7FD}
\definecolor{rule}{HTML}{C9CFF0}
\definecolor{okgreen}{HTML}{1B7A44}
\definecolor{okbg}{HTML}{F1F8F3}

\pagestyle{fancy}
\fancyhf{}
\fancyhead[L]{\footnotesize Name: \underline{\hspace{4cm}}}
\fancyhead[C]{\textbf{\color{navy}Integration Academy}\\[-2pt]{\scriptsize Calculus AB \textperiodcentered{} AP\textsuperscript{\textregistered} Exam Preparation}}
\fancyfoot[L]{\tiny\color{black!55}\parbox[t]{0.78\textwidth}{AP\textsuperscript{\textregistered} is a registered trademark of the College Board, which was not involved in the production of, and does not endorse, this material. \textcopyright{} Integration Academy.}}
\fancyfoot[R]{\scriptsize\thepage}
\renewcommand{\headrulewidth}{1.2pt}
\renewcommand{\headrule}{\hbox to\headwidth{\color{gold}\leaders\hrule height \headrulewidth\hfill}}
\renewcommand{\footrulewidth}{0pt}
\setlength{\parindent}{0pt}
\setlength{\parskip}{4pt}

% Title block: \wstitle{code}{title}{unit line}
\newcommand{\wstitle}[3]{\begin{center}{\small\bfseries\color{gold}WORKSHEET #1}\\[3pt]{\LARGE\bfseries\color{navy}#2}\\[3pt]{\small\color{black!60}#3}\end{center}\vspace{2pt}}
% Section heading
\newcommand{\sech}[1]{\par\vspace{10pt}{\large\bfseries\color{navy}#1}\par\nobreak\vspace{2pt}{\color{rule}\hrule height0.8pt}\vspace{7pt}}
% A level tag printed beside a group of questions
\newcommand{\level}[1]{\par\vspace{6pt}{\footnotesize\bfseries\color{gold}\MakeUppercase{#1}}\par\nobreak\vspace{2pt}}

\newtcolorbox{ideabox}{colframe=ink,colback=inkbg,coltitle=white,colbacktitle=ink,fonttitle=\bfseries\small,title={Key ideas},arc=2pt,boxrule=0.8pt,left=8pt,right=8pt}
\newtcolorbox{example}[1]{colframe=rule,colback=white,coltitle=navy,colbacktitle=inkbg,fonttitle=\bfseries\small,title={#1},arc=2pt,boxrule=0.8pt,borderline west={3pt}{0pt}{ink},left=10pt,right=8pt}
\newtcolorbox{frbox}[1]{colframe=gold,colback=goldbg,coltitle=white,colbacktitle=gold,fonttitle=\bfseries\small,title={#1},arc=2pt,boxrule=0.8pt,breakable,left=8pt,right=8pt}
\newtcolorbox{solbox}{colframe=okgreen,colback=okbg,coltitle=white,colbacktitle=okgreen,fonttitle=\bfseries\small,title={Worked solutions},arc=2pt,boxrule=0.8pt,breakable,left=8pt,right=8pt}
\newtcolorbox{scorebox}[1]{colframe=navy,colback=white,coltitle=white,colbacktitle=navy,fonttitle=\bfseries\small,title={#1},arc=2pt,boxrule=0.8pt,breakable,left=8pt,right=8pt}

\newcommand{\soln}{\par\smallskip\textbf{\color{ink}Solution.}\ }
\newcommand{\result}[1]{\par\smallskip\textbf{#1}}
% blank working space of a given height
\newcommand{\work}[1]{\par\vspace{#1}}
% a point value in the margin of a scoring guide
\newcommand{\pts}[1]{\hfill{\footnotesize\bfseries\color{gold}[#1]}}

% Coordinate grid for graph questions: \begin{apgraph}{xmin}{xmax}{ymin}{ymax} ... \end{apgraph}
\newenvironment{apgraph}[4]{\begin{center}\begin{tikzpicture}\begin{axis}[
  axis lines=middle, xmin=#1, xmax=#2, ymin=#3, ymax=#4,
  grid=both, grid style={black!12}, xtick distance=1, ytick distance=1,
  ticklabel style={font=\scriptsize}, xlabel={$x$}, ylabel={$y$},
  width=10cm, height=7.2cm, samples=80, clip=false,
  every axis plot/.append style={line width=1.1pt, color=ink}]}{\end{axis}\end{tikzpicture}\end{center}}
\newcommand{\closeddot}[2]{\addplot[only marks, mark=*, mark size=2.2pt, color=ink] coordinates {(#1,#2)};}
\newcommand{\opendot}[2]{\addplot[only marks, mark=*, mark size=2.2pt, color=ink, mark options={fill=white}] coordinates {(#1,#2)};}

% Multiple-choice options: four across, or stacked when the options are long.
\newcommand{\opts}[4]{\par\vspace{3pt}\begin{tabular}{@{}p{0.235\linewidth}p{0.235\linewidth}p{0.235\linewidth}p{0.235\linewidth}@{}}(A)\ #1 & (B)\ #2 & (C)\ #3 & (D)\ #4\end{tabular}\par\vspace{9pt}}
\newcommand{\optsv}[4]{\par\vspace{3pt}\begin{tabular}{@{}l@{\ \ }p{0.9\linewidth}@{}}(A) & #1\\(B) & #2\\(C) & #3\\(D) & #4\end{tabular}\par\vspace{9pt}}

\fancyhead[R]{\footnotesize\color{black!60}Worksheet 6.2}
\begin{document}

\wstitle{6.2}{Riemann Sums and the Trapezoidal Rule}{Unit 6: Integration and Accumulation of Change}

\begin{ideabox}
\begin{itemize}[leftmargin=*,itemsep=2pt,topsep=2pt]
  \item A \textbf{left}, \textbf{right} or \textbf{midpoint} Riemann sum adds rectangles whose heights are taken at the left end, right end or midpoint of each subinterval.
  \item A \textbf{trapezoidal sum} uses the average of the two end heights on each subinterval.
  \item With a table, use the table's own subintervals. The widths are often \textbf{unequal}.
  \item If $f$ is increasing, a left sum underestimates and a right sum overestimates; if $f$ is decreasing, the reverse.
\end{itemize}
\end{ideabox}

\sech{Worked examples}
\begin{example}{Example 1 \textperiodcentered{} Equal widths}
Approximate $\displaystyle\int_{0}^{3}\left(x^2+1\right)\,dx$ with left and right Riemann sums using 3 subintervals.\soln Width 1. Left: $f(0)+f(1)+f(2)=1+2+5$. Right: $f(1)+f(2)+f(3)=2+5+10$.\result{Left sum 8; right sum 17.}
\end{example}

\begin{example}{Example 2 \textperiodcentered{} Unequal widths}
Selected values of $h$ are shown.\begin{center}\renewcommand{\arraystretch}{1.25}\begin{tabular}{|c|c|c|c|c|}\hline $t$ & $0$ & $1$ & $4$ & $6$ \\ \hline $h(t)$ & $2$ & $5$ & $9$ & $7$ \\ \hline\end{tabular}\end{center}Approximate $\displaystyle\int_{0}^{6}h(t)\,dt$ with a trapezoidal sum.\soln $\dfrac12(2+5)(1)+\dfrac12(5+9)(3)+\dfrac12(9+7)(2)=3.5+21+16$.\result{$40.5$.}
\end{example}

\begin{example}{Example 3 \textperiodcentered{} Over or under}
A function is increasing on $[a,b]$. Compare its left and right Riemann sums with the integral.\soln Each left rectangle lies below the graph and each right rectangle above it.\result{Left sum $<$ integral $<$ right sum.}
\end{example}

\newpage
\sech{Practice}
Questions 5 to 7 use this table of values of a continuous function $R$.\begin{center}\renewcommand{\arraystretch}{1.25}\begin{tabular}{|c|c|c|c|c|c|}\hline $t$ & $0$ & $2$ & $3$ & $7$ & $10$ \\ \hline $R(t)$ & $6$ & $10$ & $12$ & $8$ & $2$ \\ \hline\end{tabular}\end{center}
\level{Core}
\begin{enumerate}[leftmargin=*,itemsep=2pt]
  \item Find the left Riemann sum for $f(x)=2x$ on $[0,4]$ with 4 equal subintervals.\work{2.3cm}
  \item Find the right Riemann sum for $f(x)=2x$ on $[0,4]$ with 4 equal subintervals.\work{2.3cm}
  \item Find the midpoint Riemann sum for $f(x)=x^2$ on $[0,4]$ with 2 equal subintervals.\work{2.3cm}
  \item A function is decreasing on $[a,b]$. Does a left Riemann sum overestimate or underestimate the integral?\work{2.3cm}
\end{enumerate}

\level{Exam level}
\begin{enumerate}[leftmargin=*,itemsep=2pt,start=5]
  \item Approximate $\displaystyle\int_{0}^{10}R(t)\,dt$ with a left Riemann sum using the four subintervals of the table.\work{2.8cm}
  \item Approximate the same integral with a right Riemann sum.\work{2.8cm}
  \item Approximate the same integral with a trapezoidal sum.\work{2.8cm}
  \item Find the midpoint Riemann sum for $f(x)=\dfrac1x$ on $[1,5]$ with 2 equal subintervals.\work{2.8cm}
\end{enumerate}

\begin{frbox}{Question 9 \textperiodcentered{} Free-response practice}
The speed of a train, $v(t)$ metres per second, is recorded at times $t$ seconds.\begin{center}\renewcommand{\arraystretch}{1.25}\begin{tabular}{|c|c|c|c|c|c|}\hline $t$ & $0$ & $10$ & $25$ & $40$ & $60$ \\ \hline $v(t)$ & $0$ & $12$ & $20$ & $16$ & $4$ \\ \hline\end{tabular}\end{center}
\begin{enumerate}[label=(\alph*),leftmargin=*,itemsep=2pt]
  \item Approximate $\displaystyle\int_{0}^{60}v(t)\,dt$ with a trapezoidal sum using the four subintervals of the table.\work{2.6cm}
  \item Using correct units, explain the meaning of the integral in part (a).\work{2.6cm}
  \item Use your answer to estimate the average speed of the train over the 60 seconds.\work{2.6cm}
\end{enumerate}
\end{frbox}

\newpage
\begin{solbox}
\begin{enumerate}[leftmargin=*,itemsep=4pt]
  \item $1(0+2+4+6)=12$.
  \item $1(2+4+6+8)=20$.
  \item Midpoints 1 and 3, width 2: $2(1+9)=20$.
  \item It overestimates.
  \item $(2)(6)+(1)(10)+(4)(12)+(3)(8)=94$.
  \item $(2)(10)+(1)(12)+(4)(8)+(3)(2)=70$.
  \item $\dfrac12(6+10)(2)+\dfrac12(10+12)(1)+\dfrac12(12+8)(4)+\dfrac12(8+2)(3)=16+11+40+15=82$.
  \item Midpoints 2 and 4, width 2: $2\left(\dfrac12+\dfrac14\right)=1.5$.
  \item \begin{enumerate}[label=(\alph*),leftmargin=*,itemsep=3pt]
    \item $\dfrac12(0+12)(10)+\dfrac12(12+20)(15)+\dfrac12(20+16)(15)+\dfrac12(16+4)(20)=60+240+270+200=770$.
    \item The distance the train travels, in metres, during the 60 seconds.
    \item $\dfrac{770}{60}\approx12.8$ metres per second.
  \end{enumerate}
\end{enumerate}
\end{solbox}

\end{document}
